%قبل از شروع مسئله ترجیحاً چیزی ننویسید. اما اگر خواستید بنویسید، در قالب لیستی (فقرات و شمارش) بنویسید.
%برای اضافه کردن نمره، بعدی از شروع‌مسئله، [x] را اضافه کنید و به جای x  نمره مورد نظر را قرار دهید. برای زیرمسئله‌ها هم به همین صورت است.
%\شروع‌مسئله[10]
\شروع‌مسئله
% صورت اصلی مسئله را در این‌جا بنویسید.
فرض کنید یک آرایه $n$ عضوی به نام $A$ و یک عنصر
 $x \in A$
  داریم که جایگاه دقیق آن را نمی‌دانیم. می‌خواهیم اندیس $i$ را پیدا کنیم، به طوریکه
  $A[i] = x$. 
با توجه به نکات ذکر شده به دو قسمت زیر پاسخ دهید.
  
%اگر مسئله چند قسمت داشت به این صورت عمل کنید:
\شروع{زیرمسائل}
\زیرمسئله % مسئلۀ ۱آ
فرض کنید می‌دانیم احتمال اینکه $x$ در نیمه اول آرایه باشد، $p$
 برابر این است که در نیمه دوم آرایه باشد($p \geq 1$). همچنین احتمال حضور آن در هر جایگاه برای هر دو نیمه یکسان است. پیچیدگی زمانی جستجوی خطی را در بدترین و بهترین حالت و حالت متوسط حساب کنید.



\زیرمسئله % مسئلۀ ۱ب
فرض کنید که آرایه $A$ را به صورت مرتب شده به ما داده‌اند. پیچیدگی زمانی جستجوی دودویی را در بدترین و بهترین حالت و حالت متوسط حساب کنید. الگوریتم جستجوی دودویی در زیر آورده شده است.

\lr{
	\begin{algorithm}
		\begin{algorithmic}
			\State $A \leftarrow sorted\;array$
			\State $x \leftarrow value\;to\;be\;searched$
			\State $binary\_search(1, Array\_Size)$
			\Procedure{$binary\_search(low, up)$}{}
			\State $mid = \frac{(low + up)}{2}$
			\If{$up < low$}
			\State $return \; "not \; found"$
			\EndIf
			\If{$A[mid] < x$}
			\State $return \; binary\_search(mid + 1, up)$
			\EndIf
			\If{$A[mid] > x$}
			\State $return \; binary\_search(low, mid - 1)$
			\EndIf
			\If{$A[mid] = x$}
			\State $return \; "x \; is \; at \; location" \; mid$
			\EndIf
			\EndProcedure
		\end{algorithmic}
	\end{algorithm}
}

%\زیرمسئله % مسئلۀ ۱پ
%\زیرمسئله % مسئلۀ ۱ت
\پایان{زیرمسائل}
\solution{
\begin{enumerate}
\item
در بهترین حالت، اولین خانه‌ای که بررسی می‌کنیم، همان خانه‌ای است که به دنبالش بودیم. پس پیچیدگی زمانی این حالت
$O(1)$
است. 

در بدترین حالت مجبور هستیم که تمام خانه‌ها را بررسی کنیم. پس در این حالت،‌ پیچیدگی زمانی برابر با
$O(|A|)$
است.

در حالت متوسط، اگر احتمال این که پاسخ در نیمه‌ی دوم باشد را q بنامیم، احتمال قرارداشتن پاسخ در نیمه‌ی اول pq است. حال اگر بدانیم پاسخ در یک بخش است، می‌دانیم که تمام خانه‌ها هم احتمال هستند. پس امیدریاضی تعداد خانه‌هایی که در این حالت باید طی کنیم،‌ نصف خانه‌های آن نیمه است(چرا؟). پس می‌توان نوشت:
$$pq+q = 1 \rightarrow q = \frac{1}{p+1}$$
$$ExpectedValue = pq*(\frac{|A|}{4}) + q*(\frac{|A|}{2} + \frac{|A|}{4})$$
$$ExpectedValue = \frac{p}{p+1}*\frac{|A|}{4} + \frac{1}{p+1}*\frac{3|A|}{4} = \frac{|A|(3+p)}{4(p+1)} \le \frac{|A|(3+1)}{4(0+1)} = |A| $$
هم‌چنین می‌توان گفت:
$$\frac{|A|(3+p)}{4(p+1)} \ge \frac{|A|(3+0)}{4(1+1)} = \frac{3A}{8} $$
پس اردر زمانی، همان 
$O(|A|)$
می‌ماند.
\item
بهترین حالت در این روش، مانند بهترین روش در قسمت قبل از 
$O(1)$
است.

بدترین حالت این روش، نیازمند 
$log(n)$
مرحله است.

به صورت میانگین، می‌توان هزینه‌های یافتن تمام خانه‌ها را حساب کنیم و سپس میانگین بگیریم.
$$\#Movements = 1 * 2^0 + 2 * 2^1 + 3 * 2^2 + ... + log(n) * 2 ^ {log(n) - 1}$$
یعنی نیمی از خانه‌ها هستند که حداقل به 
$log(n)$
مرحله نیاز دارند. پس میانگین از اردر 
$O(log(n))$
خواهد بود.
\end{enumerate}
}



%%%%%%%%%%%%%%%%%%%%%%%%%%%%%%%%%%%%%%%%%%%%%%%%%%%%%
% برخی دستورات مفید برای تمرین‌ها. %
%%%%%%%%%%%%%%%%%%%%%%%%%%%%%%%%%%%%%%%%%%%%%%%%%%%%%

% فرمول‌های نمونه:
%$f(x) = \Theta(x)$
%$T(x, y) \leq \log(x) + 2^y + \binom{2n}{n}$
% {\tt code\_function}


% لیست موارد بدون شماره‌گذاری:
%\شروع{فقرات}
%    \فقره مورد اول 
%        \شروع{فقرات}
%            \فقره مورد اول داخلی 
%                \شروع{فقرات}
%                    \فقره اولین مورد خیلی داخلی 
%                    \فقره دومین مورد خیلی داخلی 
%                \پایان{فقرات}
%            \فقره مورد دوم داخلی 
%        \پایان{فقرات}
%    \فقره مورد دوم 
%\پایان{فقرات}

% لیست موارد شماره‌گذاری شده:
%\شروع{شمارش}
%    \فقره یک مورد 
%    \فقره یک مورد دیگر 
%    \فقره باز هم یک مورد دیگر
%\پایان{شمارش}

% برای نوشتن یک فرمول چند خطی که یک شماره بخوره. برای حالت‌های دیگر به سایت زیر مراجعه کنید:
% https://www.overleaf.com/learn/latex/Aligning_equations_with_amsmath
%\begin{align} 
%    \begin{split}
%        (x+y)^3 &= (x+y)^2(x+y) \\
%                &= (x^2+2xy+y^2)(x+y) \\
%                &= (x^3+2x^2y+xy^2) + (x^2y+2xy^2+y^3) \\
%                &= x^3+3x^2y+3xy^2+y^3
%    \end{split}                                 
%\end{align}

% نوشتن ماتریس و آرایه دوبعدی
%\begin{align}
%    A = 
%    \begin{bmatrix}
%        A_{11} & A_{21} \\
%        A_{21} & A_{22}
%    \end{bmatrix}
%\end{align}

% اضافه کردن تصویر با شرح. برای اضافه کردن تصویر بدون شرح، فقط دستور درج‌تصویر را استفاده کنید.
%\شروع{شکل}[h!]
%\تنظیم‌ازوسط
%  \درج‌تصویر[پهنا=0.5\پهنای‌متن]{img.jpg}
%\شرح{شرح تصویر}
%\برچسب{تصویر:۱}
%\پایان{شکل}
%\پایان{زیرمسائل}

%کد زیر را ببینید.
%%\lstinputlisting{source_filename.py}
%\begin{latin}
%\begin{lstlisting}
%# Import the modules
%import sys
%import random
%
%ans = True
%
%while ans:
%question = raw_input("Ask the magic 8 ball a question: (press enter to quit) ")
%
%answers = random.randint(1,8)
%
%if question == "":
%sys.exit()
%
%elif answers == 1:
%print "It is certain"
%
%elif answers == 2:
%print "Outlook good"
%
%elif answers == 3:
%print "You may rely on it"
%
%elif answers == 4:
%print "Ask again later"
%
%elif answers == 5:
%print "Concentrate and ask again"
%
%elif answers == 6:
%print "Reply hazy, try again"
%
%elif answers == 7:
%print "My reply is no"
%
%elif answers == 8:
%print "My sources say no"
%\end{lstlisting}
%\end{latin}
